\documentclass[11pt,a4paper]{article}
\usepackage[T1]{fontenc}\usepackage{lmodern,amsmath,amssymb,amsthm,mathtools,booktabs,microtype}
\usepackage[margin=27mm]{geometry}\usepackage[colorlinks=true,linkcolor=blue,urlcolor=blue,citecolor=blue]{hyperref}
\newtheorem{theorem}{Theorem}[section]\newtheorem{lemma}[theorem]{Lemma}\newtheorem{proposition}[theorem]{Proposition}
\newcommand{\Q}{\mathbb Q}\newcommand{\Z}{\mathbb Z}\DeclareMathOperator{\Tr}{Tr}
\setlength{\emergencystretch}{2em}
\title{A 212-variable quartic with additive height comparisons}
\author{}\date{13 November 2026}
\begin{document}\maketitle
\begin{abstract}
We reduce the pinned 214-variable quartic family to 212 total variables,
under the same explicitly retained arithmetic premises. The new family
has 186 primitive quadratic equations and 2908 nonzero quartic
monomials. Replacing the comparison arguments $\lambda,T\lambda$ by
$\lambda,\lambda+T$ permits both contact calculations to use three
rather than four outputs, even after the scalar normalization of the
elliptic coordinates. A rational rescaling of two coordinates also
retains the earlier coefficient-height improvement. The sufficient
defining parameter remains ineffective. The changed arithmetic
applications, all rational boundary cases, all retained rows, and all
quartic coefficients are treated separately from the source's still
incomplete independent arithmetic audit.
\end{abstract}

\section{Statement and retained source}
Every variable count includes the distinguished input. The source is
the exact 214-coordinate, 188-row family at upstream commit
\texttt{1f7e1f1fb0438feda3c9f3da9cc16982126aeeed}. Its numeric source,
reconstructed from the pinned manuscript and checked against the
published expanded quartic, is retained in
\begin{center}\small
\nolinkurl{computation/quartic/balanced214_v24/source}.
\end{center}
The source-row SHA-256 is
\begin{center}\small
\texttt{49cf4066bf1e35e02e2156bf4bd24a21bb4c5e794f7374e056413d8306fb73c5}.
\end{center}
The coordinate list is separate from the quadratic table. Every new
row and retained row is specified by the transformation in Section~2
and by the complete output table; no generic-coordinate deletion is
being asserted.

The retained premises are the source elliptic group and index-height
statements; its scalar extraction and prepared realization; its pole
obstruction and prepared nonzero realization; the full-family
geometric/modular contact-height supplier; the aggregate \emph{sum}
height theorem; and the fixed-data simultaneous-prime statement with
arbitrary finite rational contact coefficients, finite avoidance,
conjugate reduction orders, compatible common-index congruences and
real neighborhoods. The exact supplying statements and their
applications are recorded in the pinned source's
\nolinkurl{research/quartic214/proofs/uniform-existence.md}, P03--P20.
The coefficient and field identities below do not certify those
arithmetic foundations. No additional $abc$ hypothesis is introduced.

\begin{theorem}\label{thm:main}
Under these retained inputs, there is a computable family
\[P_N\in\Z[X,Y_1,\ldots,Y_{211}],\qquad N\ge1,\]
with exact degree four, content one, and 2908 nonzero coordinate
monomials. There is one fixed
integer $N_0$ such that every $N\ge N_0$ has rational input projection
exactly $\Z$. No numerical sufficient $N_0$ is determined.

In the unbalanced presentation, put $A_N=149^{2N+1}$. Its exact
height is $\max(C_0,A_N^2+1)$, where $C_0$ is the explicit constant in
the coefficient certificate, and hence is $149^{4N+2}+1$ for $N\ge53$.
In the balanced presentation put $C=149^N$ and
\[
 K_0=31581660^{28},\qquad J_0=5898050446009651968,\qquad
 H_0=(J_0K_0)^2.
\]
Its exact height is
\[
                    \max(H_0,K_0C^2),
\]
and is $31581660^{28}149^{2N}$ for $N\ge57$.
Both presentations have the same rational input projection at each
fixed $N$. These coefficient thresholds are not defining thresholds,
and the sufficient $N_0$ is not identified with the old 214 threshold.
\end{theorem}

\section{Exact changes to the quadratic equations}
Use the pinned source coordinates before compaction. Its scalar
normalizer is $R_s=x_{118}$, its square is $w=x_{117}$, and the root
coordinates in four blocks are $C_i,E_i\in K=\Q(\sqrt2)$.
The unchanged scalar rows give
\[
 \beta=51X^2+2504,\qquad
 r_s=\beta(\beta-1)(\beta-2401)(32\beta+343)>0,
\]
\[
       (1+149^2r_s)R_s=A_Nr_s,\qquad w=R_s^2>0.
\]
These identities use only retained rows. They can also be checked
from source rows 48,49,60,63 and the input-square equation; the
intermediate pair is $(51X^2+2501)^2/4$.
The physical trace is
\begin{equation}\label{eq:scale}
 T=\frac14\Tr\sum_{i=1}^4 x([n_i]P),\qquad
 wT=\frac18\sum_{i=1}^4\left((C_{i,0}-E_{i,0})^2
                          +2(C_{i,1}-E_{i,1})^2\right).
\end{equation}
Each root is a nonzero doubled point; its abscissa exceeds $3\mp\sqrt2$
in the corresponding real embedding. Thus $T>6$.

Repurpose coordinate 4 as $t_0=\lambda$, and coordinate 90 as the
\emph{unscaled} second argument $t_1$. Replace source row 182 by
\begin{equation}\label{eq:aggregate}
 8w(t_1-t_0)=\sum_{i=1}^4
       \left((C_{i,0}-E_{i,0})^2+2(C_{i,1}-E_{i,1})^2\right).
\end{equation}
Since $w>0$, this gives $t_1=t_0+T$. The equation remains quadratic;
this is the point at which additive rather than multiplicative
comparison permits removal of the old scaling on $t_1$.

Put
\[
 \Theta_0=\{0,-1/2,-1,1/2399,-32/407\},\qquad
 \Theta=\Theta_0\cup(2-\Theta_0),
\]
and
\[
 P_5(r)=\prod_{\theta\in\Theta_0}\bigl(r-(1-\theta)^2\bigr)
       =r^5-e_1r^4+e_2r^3-e_3r^2+e_4r-e_5.
\]
The five squared distances are
$1,9/4,4,5750404/5755201,192721/165649$.
For $j=0,1$, let $q_j$ and $s_j$ be the existing weight and its
inverse; retain $q_js_j=1$. Replace the contact equations by four
quadrics with three outputs $r_j,u_j,v_j$:
\begin{align}
 r_j&=(t_j-1)^2,\qquad u_j=r_j^2,\\
 v_j&=r_ju_j-e_1(1+s_j)u_j+e_2(1+s_j)r_j,\\
 u_jv_j&=(1+s_j)(e_3u_j-e_4r_j+e_5).\label{eq:contact}
\end{align}
The first block uses old slots $(76,77,78)$ and inverse weight 80.
The second uses $(91,92,93)$ and inverse weight 95. Delete old
coordinates 119 and 120 and their defining rows 64 and 65.
The four equations replace source rows 22--25 and 35--38,
respectively. Compact the remaining coordinates and rows in increasing
old-index order. All other 177 rows are unchanged, up to this
compaction and primitive normalization. In particular no local-factor
exponent, point relation, scalar condition, or pole equation changes.
There are $214-2=212$ coordinates and $188-2=186$ rows.

\begin{lemma}\label{lem:contact}
With $qs=1$, equations \eqref{eq:contact} are equivalent to
\[
             r=(t-1)^2,\qquad (1+q)P_5(r)=r^5.
\]
The extra outputs are uniquely determined. Every solution has
$P_5(r)\ne0$. The boundary $t=1$ is included and has $q=s=-1$.
\end{lemma}
\begin{proof}
Substitute $u=r^2$ and the expression for $v$ into the last residual.
The result is exactly
\[
                         (1+s)P_5(r)-sr^5.
\]
Multiplication by $q$ gives the asserted weight equation. These
steps are reversible because $q\ne0$. If $P_5(r)=0$, then $r=0$,
but $P_5(0)=-e_5\ne0$, a contradiction. At $r=0$ the same equation
forces $s=-1$ and then $q=-1$; the outputs are $u=v=0$.
\end{proof}

Consequently $t_j$ avoids all ten contacts, including 0 and 2, but
$t_j=1$ is not discarded. The first old contact representation is
recovered, for purposes of comparison, by
\[
 \rho=1/t_0,\quad \rho^2,\rho^4,\rho^5,\quad
 G(\rho)=t_0^{-10}P_5((t_0-1)^2).
\]
The second old representation is recovered by
\[
 t_{1,\mathrm{old}}=wt_1,\quad
 r_{\mathrm{old}}=wr_1,\quad
 u_{\mathrm{old}}=wr_1^2,\quad
 v_{\mathrm{old}}=wr_1^3,\quad
 F_{\mathrm{old}}=w^2P_5(r_1).
\]
Every old contact equation is then satisfied exactly. This verifies
the same physical weight $q_j$, not a replacement chosen only to
have the right degree. The old aggregate equation is deliberately
not retained; its new arithmetic use is proved next.

\section{The changed height comparison}
Let $M(t)$ be the odd contact radical outside the source fixed
exceptional set $\Sigma$, comprising primes at most 76176 and
2029147. The source full-family height supplier, applied to the two
reflected five-contact blocks, gives a fixed $H_2>0$ with
\begin{equation}\label{eq:contactheight}
                         h(t)\le H_2M(t)^{7/6}
\end{equation}
for every rational $t$ away from the ten contacts. It includes
$t=1$, with radical one. This is the same source application as
before: the contact set and excluded primes have not changed.
The aggregate \emph{sum} theorem supplies fixed $M_1>0,M_2\ge1$,
with $M_2$ ineffective, such that
\[
       B_0^2\le M_1h(T)+M_2,\qquad B_0=\max_i|n_i|.
\]
The index tuples and trace in \eqref{eq:scale} are the identical
objects consumed by that theorem; no product-height statement is
substituted.

For the new arguments the additive height inequality gives
\begin{equation}\label{eq:addheight}
 h(T)=h(t_1-t_0)\le h(t_1)+h(t_0)+\log2.
\end{equation}
Thus whenever $M(t_j)^k\mid J$ separately for $j=0,1$, for one
nonzero integer $J$, one has
\begin{equation}\label{eq:bridge}
 B_0^2\le C|J|^{7/(6k)}+M'_2,\qquad
 C=2M_1H_2,\quad M'_2=M_2+M_1\log2.
\end{equation}
All constants are independent of $N$, the input and every witness.
Only a fixed additive constant changes. It is not necessary for the
two radicals to be coprime.

\section{Local implication and soundness}
The unchanged algebraic rows recover the same physical reference,
scalar and numerator points as in the source. The scalar normalizer
is positive independently of the changed contact and aggregate rows.
The source's forced-prime argument therefore excludes the forbidden
reference torsion and zero denominators before logarithms, exactly
as in P04 of its uniform-existence proof. The new contact lemma gives
the same $q_j=R_j^{-1}-1$ whenever $(t_j-1)\ne0$, with
$R_j=P_5((t_j-1)^2)/(t_j-1)^{10}$; no division by $t_j-1$ is used
on the boundary or unless a contact prime is being consumed.

At a consumed prime, write $v(q_j)=-\nu$, with $\nu$ positive odd.
The five unchanged physical factors are
\begin{align*}
1+q^5C_2^2+q^2D_H^2+q^4B_1^4,\\
1+q^5G^2+q^2A_3^2+q^4B_0^4,\\
1+q^5I^4+q^3A_2^2+q^2A_4^2+q^4B_3^4,\\
1+q^5C_3^2+q^2C_1^4+q^4B_4^4,\\
1+q^5C_4^2+q^2A_1^2+q^4B_2^4.
\end{align*}
Here the notation is the source notation: $D_H,A_i,B_i$ are
reciprocal abscissas, $C_i$ are differences of primary and conjugate
formal ratios, $I$ is the input comparison, and $G$ is the Lucas
residual in those ratios. Reciprocals have even valuation. The
nonconstant summands have distinct valuation classes, including the
special third factor, so negative valuations and zero tests cannot
cause cancellation. The cleared trace and pole equations force the
maximum place valuation of the factor product to exceed $-\nu/2$.
They give order at least three for $C_2,C_3,C_4,G$, order at least
one for $C_1,I$, and positive even reciprocal orders. Integrality of
one full logarithmic ratio is transferred to the other before using
relative order-four precision as an absolute congruence.

Consequently, for the same nonzero integer labels $(n,a,t,b)$,
\[
 \Delta=t^2-a(a-2)tb+(a-2)^2b^2-1,
\]
and for $J$ the reduced numerator of $n-\beta$ when nonzero,
\begin{equation}\label{eq:divs}
 M(t_j)^3\mid\Delta,\qquad M(t_j)\mid J\qquad(j=0,1).
\end{equation}
The source proof of these conclusions uses the contact equation only
at a consumed prime and uses neither $t_1/t_0=T$ nor any other
relation between the two arguments. Its hypotheses are therefore
unchanged. The same bounds $|\Delta|\le14B_0^4$ and
$|J|\le2B_0 e^{h(\beta)}$ apply.

The unchanged scalar conclusion is
\begin{equation}\label{eq:scalar}
 c=|a|\ge149^{L_N},\qquad
 h(\beta)\le2^{330770}149^{-L_N}c(1+\log c/3),
 \quad L_N=\left\lfloor\frac{9(2N-1)}4\right\rfloor+1.
\end{equation}
If $\Delta\ne0$, \eqref{eq:bridge} gives
$B_0^2\le C14^{7/18}B_0^{14/9}+M'_2$, so $B_0$ is bounded by a
fixed constant. The lower bound $B_0\ge c\ge149^{L_N}$ excludes
this case for every sufficiently large fixed $N$, uniformly in the
input and all solutions.

When $\Delta=0$, put $c=|a|$ and $w_0=\operatorname{sign}(a)(a-2)b$.
Then $t^2-ctw_0+w_0^2=1$. The norm-one unit reduction with
$\epsilon=(c+\sqrt{c^2-4})/2$ shows that $|w_0|=U_j$ and
$|t|\in\{U_{j-1},U_{j+1}\}$, where
$U_0=0,U_1=1,U_{j+1}=cU_j-U_{j-1}$.
For clarity, $\epsilon$ is the smallest norm-one unit greater than
one in $\Z[\epsilon]$: any such unit with nonzero coefficient of
$\epsilon$ differs from its inverse in absolute value by at least
$\sqrt{c^2-4}$. Reduction into $[1,\epsilon)$ proves the assertion
without assuming a maximal order. Since $w_0\ne0$ and
$|a-2|\mid w_0$, reduction of $U_j$ modulo $c-2$ or $c+2$ gives
$j\ge c-2$. Induction gives $U_j\ge(c-1)^{j-1}$. For $c\ge8$ this
implies the convenient weaker bound
\[
                        \log B_0\ge(c-5)\log(c-2).
\]

If also $J\ne0$, \eqref{eq:bridge} with $k=1$ gives, once
$B_0^2\ge2M'_2$,
\[
 \log B_0\le K+\tfrac75 h(\beta)
 \le K+\epsilon_N c(1+\log c/3),\qquad
 \epsilon_N=\tfrac75\,2^{330770}149^{-L_N},
\]
where $K$ is fixed. Choose $N$ large enough that $\epsilon_N\le1$
and $c\ge\max(e^{12},2\max(K,0)+1,\sqrt{2M'_2}+1)$.
Then
\[
(c-5)\log(c-2)-c-\tfrac c3\log c
 \ge c\bigl(\tfrac16\log c-1-\tfrac12\log2\bigr)>c/2>K,
\]
a contradiction. Hence $n=\beta$ is integral. If $X=u/v$ is
reduced, $v^2\mid51$; squarefreeness of 51 gives $v=1$.
This proves soundness with one fixed threshold. The additional
$\log2$ in \eqref{eq:addheight} changes only that ineffective
threshold, not the strict exponents.

\section{Completeness with the new arguments}
The scalar part is unchanged. Its source proof chooses a fixed
threshold $N_C$ independent of the integer input: for the positive
integer scalar product $r_s$, one has
$R_s\ge A_N/(1+149^2)$, and fixed compact real neighborhoods make
the scalar pole input lie in $(0,17)$ for every $N\ge N_C$.
Finite-place depths may depend on $(N,X)$, as before. Choose its
positive scalar label $a$ sufficiently divisible and large, then its
nonzero Lucas labels $t,b$, and $n=\beta$. Their four points and the
positive trace $T>6$ are now fixed. No subsequent step imposes a
further lower bound on $N$.

Choose fixed compact real and finite-place neighborhoods for the
common reference point and all its numerator multiples, avoiding the
source denominators. The new coefficient list in $\lambda$ is
\[
                       \Theta\ \cup\ (\Theta-T).
\]
The span of $\Theta$ is 4 and $T>6$, so there are exactly twenty
distinct coefficients. The prime-selection premise permits this
fixed list; it is not restricted to the nineteen or fewer values in
a multiplicative list with a shared zero. Choose $D_1$ so that
$D_1^2\theta$ and $D_1^2(\theta-T)$ are integers, include all fixed
avoidance and preparation primes, and take
$\lambda=u/(D_1^2w_1^2)$. The selected prime numerators are
\[
             u-\theta D_1^2w_1^2,\qquad
             u+(T-\theta)D_1^2w_1^2\quad(\theta\in\Theta).
\]
All labels, coefficients and moduli are fixed before invoking the
prime theorem. Its same conjugate annihilating orders and compatible
common-index progression give the required primary/opposite formal
groups while preserving the scalar data and real neighborhoods.
At each selected contact, $v(q)=-1$, reciprocal tests have depth at
least two, and comparison and Lucas tests have depth at least four.
Every nonconstant summand of the five factors is then positive at
the successful place. At the opposite place its inverse residue is
zero or one. The reciprocal trace is a unit of residue one or two,
exactly as in the source.

For either fixed offset $a_0\in\{0,T\}$, with $v_0=D_1^2w_1^2$,
\[
 R_j=\frac{\prod_{\theta\in\Theta}(u+(a_0-\theta)v_0)}
                {v_0^{10}(\lambda+a_0-1)^{10}}.
\]
The denominator is a square, and all numerator factors are the
selected primes, each used once in that detector. The source factor
$\gamma^2$, $\gamma=408/409$, and the squared reciprocal trace do
not change valuation parity. Hence the pole input has square reduced
denominator. At preparation primes, choose $\lambda$ to dominate
both 1 and $T$ $p$-adically; then $R_j$ is arbitrarily close to one,
$q_j$ is small and the inverse trace is close to two. The same
$\gamma$ gives local squareclasses and depths $4,2,2$ at $2,3,17$.
In the real compact neighborhoods, as $\lambda$ tends to infinity,
\[
 R_j=\prod_{\alpha}(1-\alpha/(t_j-1)^2)\longrightarrow1,
 \quad q_j>0,\quad W_j\longrightarrow1,
 \quad b_j\longrightarrow\gamma^2/4\in(0,17).
\]
An arbitrarily large fixed lower bound on the slope supplies all
these conditions. The prepared nonzero pole-realization premise now
supplies independent witnesses for the two actual pole inputs.

Finally fill the new contact outputs by \eqref{eq:contact}, and take
$t_1=\lambda+T$. The unchanged rows hold by the source physical
reconstruction, while the nine changed rows hold by construction.
This proves completeness for every integer input and every
$N\ge N_C$. Taking the maximum with the soundness threshold proves
Theorem~\ref{thm:main}. No numerical sample tuple is used to replace
this universal existence argument.

\section{Balanced coefficients and exact certificates}
Set $A_N=149C^2$, $C=149^N$. In the new unbalanced system rescale
only the coordinates with parent-214 IDs 117 and 118 by
$x_i=C y_i$. All other coordinates are fixed. Divide each transformed
row by its smallest common power of $C$ and normalize to a primitive
integer row. Every resulting row remains quadratic and is at most
linear in $C$. Since $C>0$, the change is invertible and preserves
every rational solution and the input. Unlike the earlier 214
balancing, no contact coordinate needs this rescaling.

The unbalanced quartic has 2909 parameter-separated terms; the
balanced one has 2910. Both have 2908 coordinate monomials for every
positive coefficient parameter. The only mixed balanced coefficient
profiles are $1+22201C^2$ and $1+2C^2$; no cancellation can alter
the support. An independent expansion identifies a coefficient equal
to $H_0$ and one equal to $K_0C^2$. Every other profile is either
$\ell C$, with $\ell^2\le H_0K_0$, or $a+bC^2$, with
\[
                     |a|K_0+|b|H_0\le H_0K_0.
\]
The latter bounds the profile by $\max(H_0,K_0C^2)$ on both sides
of $C^2=H_0/K_0$; the former follows by the geometric mean bound.
Thus the exact height is the asserted maximum. Exact integer
comparison gives $K_0 149^{112}<H_0\le K_0 149^{114}$, proving the
coefficient threshold 57. The unbalanced certificate similarly
identifies the unique $A_N^2+1$ profile and verifies threshold 53.

The directory \texttt{computation/quartic/additive212\_v25} contains
both complete quadratic and quartic tables, their builder, and a
checker importing no producer. The checker re-expands all 188 source
rows and matches every source coefficient, reconstructs all nine
changed equations independently, compares all 177 retained rows,
checks both physical contact normalizations and the boundary $t=1$,
and expands every new coefficient. It also verifies the two-coordinate
rescaling on every row, the exact height profiles for all positive
parameters, and deliberate rejection of a changed coefficient and a
missing row. These finite checks do not prove the source arithmetic
premises or complete an independent non-author review.

\section{Source and research status}
The new contribution is the additive comparison, its simultaneous
realization, and its compatibility with both shorter contact circuits.
The earlier attempt to delete these two outputs while preserving the
multiplicative, rescaled argument did not establish this construction.
Changing to \eqref{eq:aggregate} is what makes the unscaled second
contact circuit quadratic. This is not a fixed-witness bijection with
the original 214 system; soundness and completeness have been proved
again at the changed arguments. The coefficient balancing alone is
an exact fixed-input bijection between the two new presentations.

The 214 source is retained as a separately attributed predecessor,
not erased by its improvement. Its exact upstream commit, tables,
reviews and open audit limitations remain distinct from these new
proofs. In particular the full prime-pattern, pole-realization and
modular-height foundations are not declared exhaustively audited by
this paper. All results here have precisely the arithmetic scope
stated at the beginning.

\begin{thebibliography}{9}
\bibitem{Source} Pinned 214-variable construction and complete Version 17
manuscript, upstream commit
\texttt{1f7e1f1fb0438feda3c9f3da9cc16982126aeeed};
\nolinkurl{research/quartic214/proofs/uniform-existence.md}, P03--P20,
and \nolinkurl{manuscripts/quartic_v17.tex}, Sections A43--A47.
The numeric source and provenance are retained in the present repository.
\bibitem{Inverse} Companion paper,
\nolinkurl{papers/quartic_existence/quartic226_20261110_v21.tex},
the inverse-weight contact calculation.
\bibitem{Balance} Companion paper, \emph{Coefficient balancing for a
214-variable quartic},
\nolinkurl{papers/quartic/coefficient_balancing_20261113_v24.tex}.
\bibitem{Precision} Companion paper,
\nolinkurl{papers/local_tests/height_precision_20261113_v23.tex},
for the sharper optional exponent $2103/2048$ in place of $7/6$.
\end{thebibliography}\end{document}
